\section{Discussion}\label{sec:discussion}

\paragraph{When to use which method} FedCAST pays off whenever the byte budget is tight relative to how fast the data evolves: it spends bytes where the global objective is affected and when it is affected. When budgets are generous (above roughly 10\,\% of the bytes of sending raw data), every event-driven method converges to the centralised quality, and the choice among them matters little. k-FED is extremely frugal when clusters are well separated and each node holds few of them, exactly the regime its theory covers~\cite{dennis2021kfed}. It needs the local $k'$ (we gave it the oracle value), carries no shape information, and reacts only at its period, so it degrades on overlapping, drifting or imbalanced data. Change- and norm-triggered reporting reaches good quality but cannot be run below a data-dependent byte floor, and its thresholds must be re-tuned per dataset. FedCAST instead takes the budget itself as input.

\paragraph{What matters in FedCAST} The ablations separate the essential from the auxiliary. Prioritised partial deltas, together with an objective-linked score to rank them, deliver most of the gain: the budget is spent on the few micro-clusters that move the objective. Swap search matters under heterogeneity. The dual-ascent threshold turns the byte budget into the operating knob. The link price redistributes traffic only when budgets are loose. The novelty override turned out to be redundant with the objective-linked score. We report this negative result because it simplifies deployment: the score alone suffices to react to emerging clusters.

\paragraph{Kafka as part of the algorithm} Designing the protocol around Kafka's semantics removed whole classes of failure handling from the algorithm. Deltas can assume their predecessor because of per-key order. Retransmissions are harmless because of idempotence and sequence numbers. Nodes never resend after a server crash because of replay and a compacted changelog. Late joiners and restarted nodes obtain the model from the compacted \texttt{fsc.global}. The price is a broker, which runs in about 400\,MB.

\paragraph{Privacy} CFs reveal aggregate statistics (counts, sums, sums of squares) rather than records. They are not differentially private. Adding calibrated Gaussian noise to each delta after clipping $\lVert x\rVert$ is a natural extension~\cite{epasto2023dpstream}. The budget controller then also bounds the number of releases, and hence the privacy composition. We leave this three-way trade-off between bytes, accuracy and privacy to future work.

\paragraph{Limitations and threats to validity} (i) Scale: experiments emulate up to 50 nodes on one machine, with real Kafka for systems metrics and a validated simulator for sweeps. Wide-area deployments add effects such as cross-region replication that we do not model. (ii) Label-based metrics on real data are bounded by how well classes form k-means clusters (e.g.\ Shuttle, Letter). We therefore report the cost ratio to a centralised learner, which is the objective every method optimises. (iii) The macro-clustering assumes $k$ is known; an unknown-$k$ aggregator such as AFCL~\cite{zhang2025afcl} or density-based aggregation could replace it without changing the send policy. (iv) Data rates are replayed in stream time, so wall-clock effects of real devices (CPU contention, energy) are measured only through the Kafka experiments. (v) At the very lowest budgets (about 1\,\% of raw traffic), FedCAST cannot operate, because its bootstrap summary of $q$ micro-clusters exceeds the budget, while k-FED's $k'$ centres fit. Letting the trigger operate on compressed summaries (fewer, merged micro-clusters) would close this gap.
